\chapter{Conclusiones}

\section{Respuesta a los objetivos}
El objetivo general era determinar cuándo un conjunto pequeño de técnicas de superresolución puede
mejorar partituras digitalizadas sin convertir una reconstrucción visualmente convincente en una
afirmación falsa sobre su contenido. La comparación controlada responde que los modelos aprendidos
sí aumentan la fidelidad promedio frente a bicúbica, pero su utilidad depende de escala, severidad,
coste y revisión de los elementos musicales.

El estado del arte permitió separar interpolación, redes convolucionales, transformadores y métodos
perceptuales, y mostró que los resultados sobre imágenes naturales, documentos u OMR no demuestran
fidelidad musical. La auditoría de SMB conservó 685 elementos, identificó 260 grupos de fuente,
documentó 67 páginas no elegibles para comparación por sus anotaciones y preservó una muestra
principal intacta dentro de su única partición oficial. La degradación final definió seis condiciones normalizadas por
escala de pentagrama y la ejecución reconcilió 1.152 resultados sobre 64 obras independientes.

EDSR y SwinIR superaron a bicúbica en PSNR-Y y SSIM-Y para todas las condiciones. SwinIR destacó en
x2 limpia, x2 moderada y x4 limpia, pero EDSR igualó o superó su fidelidad en las condiciones más
difíciles y fue entre 5,2 y 8,5 veces más rápido. La revisión cualitativa halló símbolos alterados,
texto o dígitos corruptos y textura no natural en casos de ambas redes. La adaptación al dominio se
ejecutó después como estudio secundario sin reutilizar las 64 obras principales: el EDSR adaptado
mejoró al oficial en PSNR-Y y SSIM-Y en las seis condiciones sobre 20 obras nuevas, con intervalos
emparejados que no cruzaron cero. El resultado profesional es un marco condicionado en el que EDSR
preentrenado constituye la opción general por equilibrio entre fidelidad y coste, SwinIR una opción
selectiva y bicúbica una referencia rápida y transparente. La adaptación es prometedora cuando se
dispone de datos gobernados del dominio, pero no queda validada fuera de SMB ni como restauración
semántica.

En la revisión fija del modelo adaptado, cuatro de seis casos mostraron mejora sin defecto nuevo
claro y los dos casos fuertes conservaron daños ya presentes en LR. El ajuste redujo halos y
texturas y mejoró pentagramas y símbolos principales, pero no recuperó de forma fiable texto,
dígitos, alteraciones y otros elementos pequeños cuando la degradación había eliminado su
información. Por ello, su ventaja cuantitativa y visual no autoriza completar ni interpretar una
partitura sin cotejo humano.

La Tabla~\ref{tab:objective-evidence} cierra la trazabilidad entre lo prometido y lo entregado. El
ajuste fino se añadió únicamente después de asegurar el núcleo y se evaluó con un contrato separado
compatible con la integridad de la comparación principal.

\begin{table}[htbp]
  \centering
  \footnotesize
  \setlength{\tabcolsep}{3pt}
  \caption{Relación entre objetivos, evidencia y estado final. Fuente: elaboración propia.}
  \label{tab:objective-evidence}
  \begin{tabular}{p{0.31\textwidth}p{0.45\textwidth}p{0.14\textwidth}}
    \toprule
    Objetivo & Evidencia principal & Estado \\
    \midrule
    Estado del arte trazable & Síntesis crítica, fuentes primarias y criterios de selección & Cumplido \\
    Auditar y gobernar SMB & 685 registros, 260 obras, elegibilidad, licencias y hashes & Cumplido \\
    Predeclarar el experimento & Protocolo v2, muestra, pesos, métricas, semillas y reglas & Cumplido \\
    Comparar métodos emparejados & 1.152 resultados, intervalos, revisión fija y tiempos & Cumplido \\
    Evaluar adaptación acotada & 45/13/20 obras separadas; dos puntos de control y 360 resultados & Cumplido con límites \\
    Producir una guía profesional & Marco por condición, salvaguardas y casos de no uso & Cumplido con límites \\
    \bottomrule
  \end{tabular}
\end{table}

\section{Limitaciones, problemas y aprendizaje}
La principal desviación fue metodológica. El primer protocolo aplicó un desenfoque absoluto
calibrado sobre fragmentos sintéticos de pentagrama mayor. En páginas SMB completas, las condiciones
moderada y fuerte cambiaron de significado y algunas se volvieron destructivas. En lugar de
reinterpretar aquel resultado, se conservó como piloto de desarrollo, se normalizó el desenfoque
por la separación entre líneas, se seleccionaron 64 obras nuevas sin solapamiento y se repitió la
comparación completa. Esta corrección retrasó el trabajo, pero produjo una conclusión más sólida y
dejó una lección general: una degradación debe expresarse en relación con la estructura que pretende
modelar, no solo en píxeles de una muestra cómoda.

Las conclusiones se limitan a páginas SMB impresas con escala de pentagrama mensurable, a
degradaciones sintéticas y a los pesos exactos de EDSR y SwinIR. La adaptación comparte corpus y
generador de degradación entre sus roles, aunque mantenga obras disjuntas, y x4 alcanzó su mejor
validación en el límite de 2.500 pasos. Se revisaron cualitativamente seis páginas por estudio, por
lo que no se conocen tasas de fallo poblacionales. No hubo evaluación ciega por
especialistas, escaneos reales, manuscritos, medición de energía, OMR ni una tarea editorial. PSNR y
SSIM miden fidelidad de imagen, no corrección musical. La procedencia y los derechos por elemento
de SMB no están disponibles, lo que impide reproducir sus imágenes en la memoria.

El trabajo reforzó conocimientos de aprendizaje profundo, procesamiento de imagen, inferencia
estadística emparejada, diseño experimental y ejecución reproducible en entornos locales y
acelerados. También mostró el valor de conservar resultados negativos, separar exploración de
confirmación y traducir una métrica técnica en una decisión profesional con límites explícitos.

\section{Contribución y legado}
La contribución no es un nuevo modelo ni una aplicación desplegada. Es un protocolo reproducible
para estudiar superresolución de partituras, una corrección de degradación basada en escala de
pentagrama, una comparación emparejada con incertidumbre, una adaptación reproducible de EDSR, una
taxonomía de fallos musicales y un marco de decisión que incluye condiciones de no uso. Este conjunto permite a bibliotecas, archivos,
profesionales de patrimonio, intérpretes, investigadores y editoriales valorar una copia ampliada
sin confundirla con el original.

El legado técnico reside en el código Python, configuraciones congeladas, hashes de puntos de
control, manifiestos, semillas, pruebas, cuadernos y scripts de análisis. Un tercero con acceso
autorizado a SMB puede reconstruir el entorno, verificar las identidades y repetir el flujo. Los
artefactos generados conservan denominadores, resultados y limitaciones, pero no redistribuyen el
conjunto ni sus imágenes. EDSR se apoya en código MIT, SwinIR en Apache-2.0 y SMB se declara CC
BY-NC 4.0 a nivel de conjunto; cualquier reutilización debe revisar además los derechos de cada
fuente.

El efecto esperado es mejorar la toma de decisiones: favorecer copias de consulta reversibles,
conservar siempre el original y exigir cotejo humano cuando el resultado tenga consecuencias. El
protocolo también deja una base para extender la adaptación, estudiar escaneos reales u OMR sin
presentar esas líneas futuras como resultados ya validados.

\section{Relación con los estudios y competencias}
El TFG integra programación, visión por computador, aprendizaje profundo, tratamiento de datos,
estadística, gestión de proyectos y comunicación técnica. La contribución
individual comprende el análisis bibliográfico, el diseño experimental, la implementación y las
pruebas. También abarca las ejecuciones local y remota, la revisión visual, el análisis estadístico,
la interpretación profesional y la redacción. Elena Vázquez Barrachina, tutora del TFG, y Jorge
Calvo Zaragoza, director experimental, contribuyen conjuntamente a la orientación y revisión
integral del trabajo. Esta contribución no altera la autoría ni la responsabilidad individual
sobre las decisiones, la implementación y las conclusiones.

Respecto a las cinco competencias transversales vigentes, el compromiso social y medioambiental se
evidencia en el análisis de derechos, acceso al patrimonio, riesgo de información inventada y coste
computacional. La innovación y creatividad aparecen en la normalización por pentagrama, la
taxonomía de fallos y el marco condicionado de uso. El trabajo en equipo y liderazgo se concreta en
la responsabilidad individual, la coordinación explícita con la tutora y el director experimental,
y la trazabilidad de decisiones, sin atribuirles trabajo ejecutado por el estudiante. La
comunicación efectiva se desarrolla en la memoria, las figuras
y la futura defensa para un tribunal no especializado. La responsabilidad y toma de decisiones se
refleja en el bloqueo del benchmark, los protocolos previos, la corrección del error de degradación,
la apertura controlada de una adaptación secundaria, los \emph{NO-GO} restantes, la conservación
de resultados negativos y los límites explícitos.
